\documentclass[11pt,a4paper]{article}

\usepackage[utf8]{inputenc}
\usepackage[T1]{fontenc}
\usepackage[spanish]{babel}
\usepackage{amsmath}
\usepackage{amssymb}
\usepackage{booktabs}

\title{Vision Transformer aplicado a FashionMNIST:\\ Teoría e Implementación en CUDA}
\author{Rushell Zavalaga}
\date{\today}

\begin{document}

\maketitle

\begin{abstract}
Este trabajo implementa un Vision Transformer (ViT) desde cero en C++/CUDA,
sin frameworks de deep learning, y lo evalúa en la tarea de clasificación de
imágenes de FashionMNIST. Se describe la teoría detrás de la arquitectura
ViT ---patch embedding, self-attention, atención multi-cabeza y el bloque
Transformer pre-norm--- y se reporta el desempeño de un modelo compacto
(aproximadamente 11\,000 parámetros, un solo bloque Transformer) entrenado
durante 4 épocas sobre las 60\,000 imágenes de entrenamiento del dataset. El
modelo alcanza 82.45\% de exactitud en test, con una curva de pérdida y
exactitud consistentemente creciente, lo que confirma que la implementación
---incluyendo los kernels CUDA de atención, normalización y activación--- es
correcta desde el punto de vista numérico y de aprendizaje.
\end{abstract}

\section{Introducción}

Desde su introducción en el paper \emph{``An Image is Worth 16x16 Words''}
(Dosovitskiy et al., 2020), los Vision Transformers han desplazado a las
redes convolucionales (CNN) como arquitectura de referencia en visión por
computadora a gran escala. A diferencia de las CNN, que procesan la imagen
mediante filtros locales que se deslizan espacialmente, un ViT trata la
imagen como una secuencia de ``palabras visuales'' (parches) y aplica
directamente el mecanismo de self-attention introducido por Vaswani et al.
(2017) en \emph{``Attention Is All You Need''}, originalmente diseñado para
procesamiento de lenguaje natural.

El objetivo de este trabajo es doble:

\begin{enumerate}
  \item \textbf{Teórico:} explicar en detalle cómo funciona un ViT, desde la
  tokenización de la imagen en parches hasta la cabeza de clasificación,
  incluyendo las ecuaciones que gobiernan cada componente.
  \item \textbf{Práctico:} validar una implementación propia de ViT, escrita
  en C++ con kernels CUDA reales (no simulados en CPU) para cada operación
  ---proyección de parches, atención multi-cabeza, normalización de capa
  (LayerNorm), activación GELU y softmax---, entrenándola sobre FashionMNIST
  y analizando su curva de aprendizaje.
\end{enumerate}

\section{Teoría del Vision Transformer}

\subsection{Motivación: de las CNN a los Transformers en visión}

Las CNN codifican dos sesgos inductivos fuertes: \textbf{localidad} (un
píxel solo interactúa con su vecindad a través del kernel) y
\textbf{equivarianza traslacional} (el mismo filtro se aplica en toda la
imagen). Estos sesgos son útiles con poco dato, pero limitan la capacidad
del modelo para capturar relaciones \emph{globales} entre regiones lejanas
de la imagen sin apilar muchas capas.

Un Transformer, en cambio, no asume ninguna estructura espacial a priori:
cada token de la secuencia puede atender directamente a cualquier otro
token, sin importar la distancia. El costo de esa flexibilidad es que el
modelo debe \emph{aprender} la noción de posición y de vecindad desde los
datos ---de ahí la necesidad del \emph{positional embedding} (sección
\ref{sec:cls-pos})--- y que típicamente requiere más datos de entrenamiento
para generalizar bien, al carecer de los sesgos inductivos de la
convolución.

\subsection{Arquitectura general de ViT}

El flujo completo de un ViT es el siguiente:

\begin{enumerate}
  \item Imagen de entrada ($H \times W \times C$)
  \item Patch Embedding (parches $\rightarrow$ tokens)
  \item Se antepone el token \texttt{[CLS]} (token de clasificación)
  \item Se suma el Positional Embedding (posición aprendible)
  \item $N$ bloques Transformer (self-attention + MLP)
  \item LayerNorm final
  \item Se extrae el token \texttt{[CLS]}
  \item Cabeza lineal + Softmax (clasificación)
\end{enumerate}

Cada componente se detalla a continuación.

\subsection{Patch Embedding}

La imagen de entrada $x \in \mathbb{R}^{H \times W \times C}$ se divide en
$N$ parches no solapados de tamaño $P \times P$:

\[
N = \frac{H}{P} \times \frac{W}{P}
\]

Cada parche se aplana en un vector de $P^2 \cdot C$ valores y se proyecta
linealmente a la dimensión del modelo $D$ mediante una matriz de pesos
aprendible $E \in \mathbb{R}^{(P^2 C) \times D}$:

\[
z_p^{(i)} = x_p^{(i)} E, \qquad i = 1, \dots, N
\]

En la práctica, esta proyección lineal es matemáticamente equivalente a una
convolución con kernel de tamaño $P \times P$ y stride $P$ (sin
solapamiento), que es exactamente cómo está implementada en este proyecto:
se reutiliza la capa convolucional existente con
\texttt{kernel\_size = stride = patch\_size}, evitando así duplicar lógica
de proyección.

\subsection{Token de clasificación (CLS) y Positional Embedding}
\label{sec:cls-pos}

Siguiendo a BERT, se antepone a la secuencia de $N$ tokens de parche un
token aprendible adicional, el \textbf{token \texttt{[CLS]}}
($z_0 = x_{class}$), cuyo estado en la salida del último bloque Transformer
se usa como representación agregada de toda la imagen para la
clasificación:

\[
z^{(0)} = [\,x_{class};\; z_p^{(1)};\; z_p^{(2)};\; \dots;\; z_p^{(N)}\,] + E_{pos}
\]

donde $E_{pos} \in \mathbb{R}^{(N+1) \times D}$ es una tabla de embeddings
posicionales aprendibles, uno por posición en la secuencia. Sin este término
el modelo sería invariante a permutaciones de los parches, perdiendo toda
noción de estructura espacial de la imagen.

\subsection{Self-Attention y Multi-Head Attention}

El corazón del Transformer es el mecanismo de \textbf{self-attention}: cada
token de la secuencia genera tres proyecciones lineales de sí mismo
---\emph{Query} (Q), \emph{Key} (K) y \emph{Value} (V)--- mediante matrices
de pesos aprendibles $W^Q, W^K, W^V \in \mathbb{R}^{D \times D}$:

\[
Q = ZW^Q, \qquad K = ZW^K, \qquad V = ZW^V
\]

La atención de un token hacia todos los demás se calcula comparando su
Query con todas las Keys de la secuencia, normalizando con softmax y usando
el resultado como pesos para combinar los Values:

\[
\mathrm{Attention}(Q, K, V) = \mathrm{softmax}\!\left(\frac{QK^\top}{\sqrt{d_k}}\right) V
\]

El factor $1/\sqrt{d_k}$ (donde $d_k = D / h$ es la dimensión por cabeza)
evita que los productos punto crezcan en magnitud con $d_k$ y saturen el
softmax, lo que produciría gradientes cercanos a cero.

En vez de una sola atención, ViT usa \textbf{atención multi-cabeza}: se
dividen Q, K, V en $h$ subespacios de dimensión $d_k = D/h$, se calcula la
atención en paralelo en cada uno, y los resultados se concatenan y se
proyectan de vuelta a $D$ con una matriz de salida $W^O$:

\[
\mathrm{MultiHead}(Q,K,V) = \mathrm{Concat}(\mathrm{head}_1, \dots, \mathrm{head}_h)\,W^O,
\qquad \mathrm{head}_i = \mathrm{Attention}(QW_i^Q, KW_i^K, VW_i^V)
\]

Tener varias cabezas permite que el modelo atienda simultáneamente a
diferentes tipos de relaciones (por ejemplo, una cabeza podría capturar
bordes cercanos y otra, simetrías globales de la prenda), en lugar de
promediar todo en una sola distribución de atención.

\subsection{Bloque Transformer (pre-norm)}

Cada bloque Transformer combina atención y una red feed-forward (MLP), cada
una envuelta en una conexión residual y precedida por \textbf{LayerNorm}
(variante \emph{pre-norm}, más estable numéricamente que la variante
\emph{post-norm} del Transformer original):

\begin{align*}
x' &= x + \mathrm{MultiHead}(\mathrm{LN}(x)) \\
x'' &= x' + \mathrm{MLP}(\mathrm{LN}(x'))
\end{align*}

\textbf{LayerNorm} normaliza cada token de forma independiente (a
diferencia de BatchNorm, que normaliza a través del batch), a media 0 y
varianza 1, y luego aplica una escala y desplazamiento aprendibles
$\gamma, \beta$ por canal:

\[
\mathrm{LN}(x)_j = \gamma_j \cdot \frac{x_j - \mu}{\sqrt{\sigma^2 + \epsilon}} + \beta_j,
\qquad \mu = \frac{1}{D}\sum_{j=1}^{D} x_j, \quad \sigma^2 = \frac{1}{D}\sum_{j=1}^{D}(x_j-\mu)^2
\]

El \textbf{MLP} consiste en dos capas lineales con una activación no lineal
\textbf{GELU} (Gaussian Error Linear Unit) entre medio:

\[
\mathrm{MLP}(x) = W_2\,\mathrm{GELU}(W_1 x + b_1) + b_2, \qquad
\mathrm{GELU}(x) = x \cdot \Phi(x) = \tfrac{1}{2}x\left(1 + \mathrm{erf}\!\left(\tfrac{x}{\sqrt{2}}\right)\right)
\]

GELU se prefiere sobre ReLU en Transformers porque es suave (derivable en
todo punto) y pondera la entrada por su probabilidad bajo una gaussiana
estándar, en vez de recortarla abruptamente en 0.

Las conexiones residuales ($x + \cdot$) son las que permiten apilar muchos
bloques sin que el gradiente se desvanezca: durante el backward, cada suma
reparte el gradiente sin modificarlo hacia ambas ramas.

\subsection{Cabeza de clasificación}

Tras pasar por los $N$ bloques Transformer y una última LayerNorm, se
extrae únicamente el token \texttt{[CLS]} (posición 0 de la secuencia)
---que a esta altura ha podido atender a todos los parches y agregado
información global de la imagen--- y se proyecta con una capa lineal a la
cantidad de clases, seguida de softmax:

\[
\hat{y} = \mathrm{softmax}\left(W_{head}\, z_L^{(0)} + b_{head}\right)
\]

El entrenamiento minimiza la pérdida de \textbf{entropía cruzada
categórica} entre $\hat{y}$ y la etiqueta real (codificada one-hot):

\[
\mathcal{L} = -\sum_{c=1}^{C} y_c \log(\hat{y}_c)
\]

\subsection{Complejidad computacional}

El costo dominante de self-attention es $O(N^2 D)$ (por el producto
$QK^\top$, de tamaño $N \times N$), frente a $O(N D^2)$ del MLP. Para
secuencias cortas ---como en este trabajo, donde $N=16$ parches más el
token CLS dan una secuencia de longitud 17--- el término $N^2 D$ es
manejable, pero es la razón por la que ViT escala peor que una CNN a
resoluciones muy altas: duplicar $H$ y $W$ cuadruplica $N$ y por lo tanto
cuadruplica (al cuadrado) el costo de atención.

\section{Caso de estudio: FashionMNIST}

\subsection{Problema}

FashionMNIST (Xiao et al., 2017) es un dataset de clasificación de
imágenes publicado por Zalando Research como reemplazo directo (mismo
formato, mismas dimensiones) del clásico MNIST de dígitos manuscritos, pero
con una tarea notablemente más difícil: 70\,000 imágenes en escala de
grises de $28 \times 28$ píxeles, repartidas en 60\,000 de entrenamiento y
10\,000 de prueba, distribuidas uniformemente en 10 clases de prendas de
vestir (camiseta, pantalón, pulóver, vestido, abrigo, sandalia, camisa,
zapatilla, bolso, botín).

La dificultad frente a MNIST radica en que las clases tienen mucha más
variabilidad intra-clase (una ``camisa'' y un ``abrigo'' pueden compartir
siluetas muy parecidas) y menos estructura de trazo simple que un dígito,
lo que la convierte en un banco de pruebas más exigente y realista para
validar si una arquitectura ---en este caso, un ViT implementado desde
cero--- realmente aprende representaciones útiles y no solo memoriza
patrones triviales de trazo.

\subsection{Marco teórico: instancia concreta del modelo}

Se instanció la arquitectura de la sección 2 con los siguientes
hiperparámetros:

\begin{table}[h]
\centering
\begin{tabular}{ll}
\toprule
\textbf{Hiperparámetro} & \textbf{Valor} \\
\midrule
Resolución de imagen & $28 \times 28 \times 1$ \\
Tamaño de parche $P$ & 7 \\
Parches por imagen $N$ & $4 \times 4 = 16$ \\
Longitud de secuencia ($N+1$ con CLS) & 17 \\
Dimensión de embedding $D$ & 32 \\
Cabezas de atención $h$ & 4 ($d_k = 8$ por cabeza) \\
Bloques Transformer & 1 \\
Dimensión oculta del MLP & 64 \\
Parámetros totales (aprox.) & $\sim$11\,100 \\
Optimizador & SGD, sin momentum \\
Tasa de aprendizaje & 0.01 \\
Tamaño de batch & 8 \\
Función de pérdida & Entropía cruzada categórica \\
\bottomrule
\end{tabular}
\caption{Hiperparámetros del modelo ViT usado en los experimentos.}
\label{tab:hparams}
\end{table}

Es un modelo deliberadamente pequeño (un solo bloque Transformer,
aproximadamente 11\,000 parámetros) frente a los ViT usados en la
literatura (decenas de millones de parámetros, 12+ bloques), tanto por
restricciones de tiempo de cómputo como para mantener la implementación
---y su verificación numérica vía \emph{gradient checking}--- manejable.

\textbf{Implementación.} Todas las operaciones del forward y backward
---proyección de parches (equivalente a convolución), separación/combinación
de cabezas de atención, el producto batched $QK^\top$ y
$\mathrm{Attn}\cdot V$, softmax y su derivada, LayerNorm (forward y
backward, con acumulación atómica de $\partial\gamma,\partial\beta$ sobre
todas las filas), GELU, y la inserción/extracción del token CLS--- están
implementadas como kernels CUDA reales, no como bucles en CPU con
transferencias de memoria disfrazadas de ``GPU''. Esto se verificó de dos
formas: (1) \emph{gradient checking} numérico sobre cada capa (diferencias
finitas vs.\ backward analítico, error máximo del orden de
$10^{-3}$--$10^{-4}$), y (2) comparando el tiempo de entrenamiento antes y
después de mover cada operación a un kernel: el tiempo por época bajó de
aproximadamente 4 minutos (con varias operaciones aún corriendo en CPU) a
un promedio de $\sim$161 segundos por época en la corrida reportada en la
sección \ref{sec:resultados}.

\subsection{Resultados}
\label{sec:resultados}

Se entrenó el modelo durante 4 épocas completas sobre las 60\,000 imágenes
de entrenamiento (con \emph{shuffle} de índices en cada época), evaluando
pérdida y exactitud sobre el set de entrenamiento completo y el set de test
completo (10\,000 imágenes) al final de cada época:

\begin{table}[h]
\centering
\begin{tabular}{cccccc}
\toprule
\textbf{Época} & \textbf{Train Loss} & \textbf{Train Acc} & \textbf{Test Loss} & \textbf{Test Acc} & \textbf{Tiempo (s)} \\
\midrule
1 & 0.5548 & 0.7979 & 0.5958 & 0.7825 & 153.3 \\
2 & 0.4832 & 0.8246 & 0.5274 & 0.8077 & 168.6 \\
3 & 0.4493 & 0.8358 & 0.5009 & 0.8182 & 159.8 \\
4 & 0.4238 & 0.8427 & 0.4752 & 0.8245 & 164.6 \\
\bottomrule
\end{tabular}
\caption{Pérdida y exactitud por época sobre FashionMNIST (train y test completos).}
\label{tab:resultados}
\end{table}

\textbf{Análisis.} La pérdida decrece de forma monótona tanto en
entrenamiento como en test, y la exactitud en test sube consistentemente de
78.25\% a 82.45\% en solo 4 épocas, sin señales de sobreajuste (la brecha
entre train acc y test acc se mantiene acotada, entre 1.5 y 1.8 puntos
porcentuales durante todo el entrenamiento). Esto indica que:

\begin{itemize}
  \item El mecanismo de self-attention implementado efectivamente propaga
  gradientes útiles de principio a fin de la red (patch embedding
  $\rightarrow$ CLS $\rightarrow$ atención $\rightarrow$ MLP
  $\rightarrow$ cabeza de clasificación), lo cual no es evidente a priori
  en una implementación construida desde cero sin un framework que valide
  automáticamente las formas y gradientes.
  \item El modelo aún está \emph{subajustando} (no ha convergido): con un
  solo bloque Transformer y $D=32$, la capacidad es baja comparada con los
  90--91\% de exactitud que reportan CNN convencionales de mayor capacidad
  sobre FashionMNIST, o con ViT de mayor escala. La tendencia estrictamente
  creciente en las 4 épocas sugiere que más épocas de entrenamiento
  ---o un modelo con más bloques/dimensión--- seguirían mejorando la
  exactitud antes de estancarse.
  \item El tiempo por época ($\sim$161\,s en promedio, con las 60\,000
  imágenes en batches de 8) es consistente entre épocas, lo que confirma
  que no hay degradación de memoria ni fugas de recursos en los kernels
  CUDA a lo largo del entrenamiento.
\end{itemize}

\subsection{Conclusiones}

La implementación desde cero de un Vision Transformer en C++/CUDA, sin
depender de frameworks como PyTorch o TensorFlow, logra aprender
representaciones útiles sobre FashionMNIST: con un modelo compacto de
aproximadamente 11\,000 parámetros y solo 4 épocas de entrenamiento, se
alcanza 82.45\% de exactitud en test, con una curva de aprendizaje estable
y sin sobreajuste. Esto valida tanto la correctitud matemática de cada
componente ---patch embedding, atención multi-cabeza, LayerNorm, GELU,
softmax y sus respectivos backwards--- como la correctitud de su portado a
kernels CUDA reales, que además reduce el tiempo de entrenamiento por época
en aproximadamente un 33\% respecto a una versión con partes del cómputo
aún en CPU. Los resultados sugieren que, con más capacidad (más bloques
Transformer, mayor dimensión de embedding) y más épocas de entrenamiento,
el modelo tiene margen para acercarse al desempeño de arquitecturas
convolucionales de referencia sobre este mismo dataset.

\section{Conclusiones generales}

Este trabajo cubrió, de extremo a extremo, la teoría y la implementación
práctica de un Vision Transformer: desde la tokenización de una imagen en
parches y su proyección lineal, pasando por el mecanismo de self-attention
y su generalización multi-cabeza, hasta el bloque Transformer completo con
normalización, MLP y conexiones residuales. La validación experimental
sobre FashionMNIST ---un dataset deliberadamente más difícil que MNIST---
confirma que la arquitectura, implementada íntegramente con kernels CUDA
propios, aprende de forma estable y consistente, sentando una base
funcional sobre la cual escalar el modelo (más bloques, mayor dimensión,
más épocas) en trabajo futuro.

\begin{thebibliography}{9}

\bibitem{dosovitskiy2020}
Dosovitskiy, A., Beyer, L., Kolesnikov, A., et al. (2020).
\emph{An Image is Worth 16x16 Words: Transformers for Image Recognition at Scale.}
arXiv:2010.11929.

\bibitem{vaswani2017}
Vaswani, A., Shazeer, N., Parmar, N., et al. (2017).
\emph{Attention Is All You Need.} NeurIPS 2017.

\bibitem{xiao2017}
Xiao, H., Rasul, K., \& Vollgraf, R. (2017).
\emph{Fashion-MNIST: a Novel Image Dataset for Benchmarking Machine Learning Algorithms.}
arXiv:1708.07747.

\bibitem{ba2016}
Ba, J. L., Kiros, J. R., \& Hinton, G. E. (2016).
\emph{Layer Normalization.} arXiv:1607.06450.

\bibitem{hendrycks2016}
Hendrycks, D., \& Gimpel, K. (2016).
\emph{Gaussian Error Linear Units (GELUs).} arXiv:1606.08415.

\end{thebibliography}

\end{document}
